
\section{10. Limitations}

\subsection{10.1 Task Coverage}

Only SST-2 is evaluated. A3 may behave differently on:

\begin{itemize}
\item Generation tasks (summarization, translation)
\end{itemize}
\begin{itemize}
\item Multi-class classification (MNLI, AG News)
\end{itemize}
\begin{itemize}
\item Token-level tasks (NER, POS)
\end{itemize}
\begin{itemize}
\item Long-context tasks (>512 tokens)
\end{itemize}

\subsection{10.2 Model Coverage}

We test on models with ≤124M parameters. Not tested:

\begin{itemize}
\item Large models (>1B): SVD time grows quadratically
\end{itemize}
\begin{itemize}
\item Very deep models (>24 layers): compounding effect
\end{itemize}
\begin{itemize}
\item Mixture-of-Experts: not tested
\end{itemize}
\begin{itemize}
\item Vision transformers: not tested
\end{itemize}
\begin{itemize}
\item Multimodal: not tested
\end{itemize}

\subsection{10.3 Hardware}

All experiments on RTX 4070 Laptop (8 GB). Behavior on other hardware:

\begin{itemize}
\item A100/H100: SVD faster, benefit larger
\end{itemize}
\begin{itemize}
\item CPU: matmul-bound; A3 might help
\end{itemize}
\begin{itemize}
\item Edge: not tested
\end{itemize}

\subsection{10.4 Calibration Dependency}

A3 requires a calibration pass. If calibration data is unavailable:

\begin{itemize}
\item Fisher weighting degrades to unweighted SVD
\end{itemize}
\begin{itemize}
\item LayerNorm recalibration is skipped
\end{itemize}
\begin{itemize}
\item Accuracy may drop by up to 5\%
\end{itemize}

This is a real limitation for models whose deployment environment differs 
from training data.

\subsection{10.5 Inference Latency}

Factorized layers are slower at inference:

\begin{itemize}
\item Two matmuls instead of one
\end{itemize}
\begin{itemize}
\item Intermediate rank-$r$ tensor must be materialized
\end{itemize}
\begin{itemize}
\item On GPU: ~1.3× slower
\end{itemize}
\begin{itemize}
\item On CPU: potentially faster
\end{itemize}

This is a \textbf{trade-off}: smaller size for slower latency.

\subsection{10.6 Single-Seed Results}

All main results are from one seed. With N=872 validation examples, SE is 
~1.69\%. Differences below 3.4\% (2×SE) may be noise.

\subsection{10.7 No Theoretical Guarantee}

We provide no formal bound on downstream accuracy loss. The Fisher proxy 
$\mathbb{E}[x^2]$ is not the true Fisher; our approximation error is untested.

\subsection{10.8 Rank Selection}

Our adaptive rank escalation increases rank until error is met. This is 
greedy and may not be optimal. A learned or Bayesian approach could be better.

\subsection{10.9 LayerNorm Recalibration Heuristic}

Our $\gamma \leftarrow \gamma \cdot \sigma_\text{target} / \sigma_\text{current}$ 
is a heuristic. It does not account for:

\begin{itemize}
\item Nonlinear interactions between layers
\end{itemize}
\begin{itemize}
\item Changes in higher moments
\end{itemize}
\begin{itemize}
\item Attention-specific distribution shifts
\end{itemize}

\subsection{10.10 No Comparison with Some Baselines}

We do not directly compare with:

\begin{itemize}
\item DistilBERT / TinyBERT (distillation)
\end{itemize}
\begin{itemize}
\item Performer / Linformer (attention replacement)
\end{itemize}
\begin{itemize}
\item SparseGPT / Wanda (pruning)
\end{itemize}

This limits the paper's external validity.

\subsection{10.11 Failure Modes Not Fully Diagnosed}

{{UNRESOLVED_FAILURES}}

\subsection{10.12 Code Quality}

Our implementation is research-grade:

\begin{itemize}
\item No unit tests for edge cases
\end{itemize}
\begin{itemize}
\item Some error paths use broad exception handling
\end{itemize}
\begin{itemize}
\item Numerical stability could be improved in some paths
\end{itemize}

\subsection{10.13 Ethical Considerations}

Compression methods enable broader deployment. This is a \textbf{dual-use} 
capability. See Appendix E.